%! Author = sbbfti
%! Date = 10/06/2020

%% Use the option review to obtain double line spacing
%% \documentclass[preprint,review,12pt]{elsarticle}

%% Use the options 1p,twocolumn; 3p; 3p,twocolumn; 5p; or 5p,twocolumn
%% for a journal layout:
% \documentclass[final,1p,times,nonatbib]{elsarticle-deutsch}
% \documentclass[final,1p,times,twocolumn,nonatbib]{elsarticle-deutsch}
%% \documentclass[final,3p,times,nonatbib]{elsarticle-deutsch}
% \documentclass[final,3p,times,twocolumn,nonatbib]{elsarticle-deutsch}
%% \documentclass[final,5p,times,nonatbib]{elsarticle-deutsch}
\documentclass[final,5p,times,twocolumn,nonatbib]{elsarticle-deutsch}

% \documentclass[ngerman,final,3p,times,nonatbib]{elsarticle-deutsch}

\usepackage{setspace}
\usepackage{mypreamble}

\usepackage[nonumberlist,nogroupskip]{glossaries}
% \makenoidxglossaries

\newglossaryentry{7730}
{
name={ISO~7730:2005},
description={ISO~7730:2005 is a thermal comfort standard developed by ISO}
}

\newglossaryentry{55}
{
name={ASHRAE~55--2020},
description={ASHRAE~55--2020 is a thermal comfort standard developed by ANSI and ASHRAE}
}

\newacronym{iupac}{IUPAC}{International Union of Pure and Applied Chemistry}

\newacronym{ptfe}{PTFE}{Polytetrafluorethylen}
\newacronym{1k}{1K-Klebstoff}{Einkomponenten-Klebstoff}
\newacronym{2k}{2K-Klebstoff}{Zwei\-kom\-ponenten-Klebstoff}
\newacronym{btex}{BTEX}{Benzol, Toluol, Ethylbenzol und Xylol}
\newacronym{voc}{VOC}{Volatile Organic Compounds}
\newacronym{who}{WHO}{World Health Organisation}
\newacronym{m-cte}{$\mathrm{\mu}$-CTE}{Micro-Chamber/Thermal Extractor}
\newacronym{dvdt}{$\dot{V}$}{Volumenstrom}
\newacronym{vd}{$V_D$}{Durchbruchvolumen}
\newacronym{td}{TD}{Thermodesorber}
\newacronym{gc}{GC}{Gaschromatograph}
\newacronym{ms}{MS}{Massenspektrometer}

% \newglossaryentry{u-cte}
% {
%     name={$\mu$-CTE},
%     description={Micro-Chamber/Thermal Extractor}
% }


\begin{document}
    
    \onehalfspacing

    \begin{frontmatter}

    \title{Titel des Papers}

    \author[label1]{Duanghathai Hartwig}
    \ead{duanghathai.hartwig@stud.hs-emden-leer.de}
    % \author[label2]{Name Surname}

    \affiliation[label1]{organization={University of Applied Sciemces Emden-Leer},
             addressline={Constantia},
             city={Emden},
             postcode={26723},
             state={Niedersachsen},
             country={Germany}}

    \cortext[mycorrespondingauthor]{Diese Arbeit wurde im Rahmen des Moduls \textit{Modul} an der Hochschule Emden-Leer verfasst.}

    
    \begin{abstract}
        Zusammenfassung...
    \end{abstract}

    % \begin{graphicalabstract}
        % \includegraphics[width=\linewidth]{figures/example_clima}
    % \end{graphicalabstract}

    % Research highlights
    % \begin{highlights}
    %     \item 
    %     \item 
    %     \item
    % \end{highlights}

    % \begin{keyword}
        % Phosgensynthese \sep Matlab/Simulink \sep Reaktionskinetik \sep Prozessmodellierung \sep Festbettreaktor
    % \end{keyword}

\end{frontmatter}

    % \tableofcontents
    % \cleardoublepage
    
    
    %% The lineno packages adds line numbers. Start line numbering with
    %% \begin{linenumbers}, end it with \end{linenumbers}. Or switch it on
    %% for the whole article with \linenumbers.
    % \usepackage{lineno}
    
    % \linenumbers
    
    \acresetall
    
    \section{Einleitung}\label{sec:einleitung}
\blindtext
\begin{figure}[H]
    \begin{adjustbox}{width=\columnwidth}
        \schemestart
        \chemname{
\chemfig{
               % 2
                  (
    -[:150,,,,drhs]% 1
                  )
         -[:30]% 3
                  (
            =[:330]O% 4
                  )
     -[:90,,,1]OH% 5
}
}{acrylic acid}
        \+
        \chemname{
\chemfig{
               % 1
        -[:150]% 2
    -[:210,,,2]HO% 3
}
}{ethanol}
        \arrow
        \chemname{
\chemfig{
                   % 1
            -[:150]% 2
            -[:210]O% 3
            -[:150]% 4
                      (
                 =[:90]O% 5
                      )
            -[:210]% 6
    -[:150,,,,drhs]% 7
}
}{acrylic acid ethyl ester}
        \+
        \chemname{
\chemfig{
    H_2O% 1
}
}{water}
        \schemestop
    \end{adjustbox}
\end{figure}
This
    
    % \newpage
    \section{Theorie}\label{sec:theorie}
\subsection{Adsorbtionsmittel zur Anreicherung der Analyten}
Voraussetzung für die Quantifizierbarkeit der emittierten Verbindungen ist, diese in erfassbaren Konzentrationsbereichen so verlustfrei wie möglich anzureichern.
Zur Verfügung stehen dafür die physikalische und chemische Adsorbtion in Flüssigkeiten, die Kondensation bei tiefen Temperaturen, sowie die Adsorbtion an Festkörpern.
Letztere hat sich für die Spurenanalytik von organischen Luftinhaltsstoffen aufgrund des heutzutage bestehenden großen Angebots an festen Sorbensmaterialien etabliert. \par
Genannte feste Sorbenzien lassen sich in anorganische Adsorbtionsmittel, Adsorbtionsmittel auf Kohlenstoffbasis und poröse Polymere einordnen.
Unerlässliche Anforderungen sind dabei, dass das gewählte Adsorbtionsmittel bei Umgebungstemperatur ein großes \ac{vd} für organische Verbindungen wie z.B. \gls{btex} aufweist, die Analyten nach der Adsorbtion wieder vollständig zu eluieren vermag und während beider Prozessbedingungen eine hohe physikalische und chemische Stabilität aufweist. \par
Das \ac{vd}, angegeben in Liter Gas pro Gramm Adsorbens, ist dabei eine auf die Adsorbensmasse bezogene Größe, welche die Gasmenge angibt die das Adsorbtionsröhrchen durchfließen darf bevor eine Spezies eluiert. 
\ac{vd} ist stark temperaturabhängig und individuell für jede adsorbierte Spezies \cite{figge_adsorptionsmittel_1987}. \par
Bei der Wahl des Adsorbens kommt es dabei vor allem auf die Flüchtigkeit des Analyten an.
So sind für sehr flüchtige Verbindungen starke, für halbflüchtige Verbindungen eher schwache Adsorbens Materialien geeignet.

    
    \section{Material und Methoden}\label{sec:materialundmethoden}


    \section{Ergebnisse}\label{sec:ergebnisse}


    \section{Diskussion}\label{sec:diskussion}
    
    % \clearpage
    
    %! Author = federico
%! Date = 10/06/2020

\section*{Glossar}
\renewcommand{\baselinestretch}{0.75}\normalsize
\renewcommand{\aclabelfont}[1]{\textsc{\acsfont{#1}}}
\begin{acronym}[longest]
    \acro{dvdt}[$\dot{V}$]{Volumenstrom}
    \acro{vd}[$V_D$]{Durchbruchvolumen}
\end{acronym}
\renewcommand{\baselinestretch}{1}\normalsize


    % \printnoidxglossaries

    \printbibliography

    % \clearpage

    % \appendix

\section{Appendix}\label{sec:appendix}


\end{document}